\chapter{Conclusion partielle}

La partie étudiée dans ce document montre qu'un système de détection de signaux
peut constituer, à lui seul, un sous-ensemble cohérent du projet global. Sa
valeur ne réside pas uniquement dans la présence d'indicateurs techniques, mais
dans l'organisation d'une chaîne méthodologique complète, allant de la génération
structurée des variantes jusqu'à leur restitution dans un environnement de
supervision.

Le coeur de cette contribution se situe dans le pipeline A-G. La couche A fournit
un univers candidat calibré par horizon ; la couche B impose une évaluation hors
échantillon rigoureuse ; la couche C transforme les performances observées en un
score de robustesse multi-critère ; les couches D et E éliminent successivement
les variantes insuffisantes puis redondantes ; enfin, les couches F et G
produisent un signal courant agrégé, interprétable et pondéré par la fiabilité
historique. Cette chaîne permet d'éviter deux écueils fréquents : la sélection
naïve du meilleur indicateur et la restitution d'un score non explicable.

L'étude des interfaces montre ensuite que la cohérence du projet ne s'arrête pas
au coeur quantitatif. La page Signaux reprend la structure logique du pipeline en
la traduisant en niveaux de lecture successifs, tandis que le tableau de bord
convertit ces résultats en outil de priorisation par titre et par horizon. La
qualité de restitution n'est donc pas ajoutée après coup ; elle prolonge la
philosophie méthodologique du moteur.

Au terme de cette analyse, la chaîne
\textit{données validées \rightarrow pipeline de signal \rightarrow page
Signaux \rightarrow tableau de bord}
peut être considérée comme une base décisionnelle crédible, traçable et
professionnellement exploitable. Elle prépare naturellement les étapes aval du
projet, notamment la transformation des signaux en règles de stratégie puis leur
validation complète, mais elle possède déjà sa cohérence propre et sa valeur
ingénierique autonome.
